\documentclass[12pt,aspectratio=1610]{beamer}
\geometry{paperwidth=20cm,paperheight=12.5cm}
\usepackage[italian,provide=*]{babel}
\usepackage{lmodern,fontspec}
\setsansfont{NimbusSans-Regular.otf}[Path=sorgenti/font/,BoldFont=NimbusSans-Bold.otf,ItalicFont=NimbusSans-Italic.otf]
\setmainfont{NimbusSans-Regular.otf}[Path=sorgenti/font/,BoldFont=NimbusSans-Bold.otf,ItalicFont=NimbusSans-Italic.otf]
\setmonofont{lmmono10-regular.otf}[Path=sorgenti/font/,Scale=.88]
\usefonttheme{professionalfonts}
\usepackage{amsmath,amssymb,booktabs,array,tabularx,graphicx,tikz,listings,bookmark}
\usetikzlibrary{calc,arrows.meta,positioning}
\definecolor{Paper}{HTML}{F8F6F1}
\definecolor{Ink}{HTML}{253344}\definecolor{Navy}{HTML}{15283E}
\definecolor{Quantum}{HTML}{5146A5}\definecolor{Classical}{HTML}{087C74}
\definecolor{Measure}{HTML}{A05B10}\definecolor{Rose}{HTML}{9B4564}
\definecolor{Muted}{HTML}{657180}\definecolor{Rule}{HTML}{D6DAD9}
\definecolor{Wash}{HTML}{EAE7F6}\definecolor{Night}{HTML}{112236}
\definecolor{Lilac}{HTML}{ABB5FF}\definecolor{Mint}{HTML}{72D8C8}\definecolor{Gold}{HTML}{EEC184}
\setbeamercolor{background canvas}{bg=Paper}\setbeamercolor{normal text}{fg=Ink}
\setbeamertemplate{navigation symbols}{}\setbeamertemplate{headline}{}\setbeamertemplate{footline}{}
\setbeamersize{text margin left=0pt,text margin right=0pt}\setbeameroption{hide notes}
\newsavebox{\BodyBox}\newsavebox{\SideBox}\newsavebox{\RecallBox}\newsavebox{\TitleBox}
\newcommand{\MeasureBox}[2]{\typeout{GEOBOX:#1:\the\wd#2:\the\ht#2:\the\dp#2}}
\lstset{language=R,basicstyle=\ttfamily\fontsize{10.2}{13.4}\selectfont,
columns=fullflexible,keepspaces=true,showstringspaces=false,breaklines=true,breakatwhitespace=false,
commentstyle=\color{Muted},keywordstyle=\color{Classical},stringstyle=\color{Quantum},
aboveskip=0pt,belowskip=0pt,xleftmargin=0pt,xrightmargin=0pt,tabsize=2}
\hypersetup{colorlinks=true,linkcolor=Classical,urlcolor=Classical,pdfpagemode=UseOutlines,pdfstartview=Fit,bookmarksopen=true,bookmarksopenlevel=0}
\bookmarksetup{open,openlevel=0,depth=2}
\tikzset{arr/.style={-{Stealth[length=2mm]},line width=.8pt,draw=Classical}}
\hypersetup{pdftitle={Campionamento, stima e incertezza},pdfauthor={},pdfsubject={Corso di statistica e data science per Geologia, primo anno; modulo 05}}
\usetikzlibrary{decorations.pathreplacing,fit,backgrounds,matrix}\hypersetup{pdftitle={P01 · Geometria dei pori e permeabilità},pdfsubject={Revisione integrale R6}}\begin{document}
\begin{frame}[plain]\hypertarget{slide1}{}\begin{tikzpicture}[remember picture,overlay]\fill[Night] (current page.south west) rectangle (current page.north east);\draw[Mint,line width=2pt] ([xshift=1.05cm,yshift=-1.35cm]current page.north west)--([xshift=3cm,yshift=-1.35cm]current page.north west);\node[anchor=north west,inner sep=0] at ([xshift=1.05cm,yshift=-1.7cm]current page.north west){\color{Mint}\bfseries\fontsize{11}{14}\selectfont P01 · REVISIONE R6};
\node[anchor=north west,inner sep=0] at ([xshift=1.05cm,yshift=-2.65cm]current page.north west){\begin{minipage}{15.8cm}\raggedright\color{white}\bfseries\fontsize{31}{38}\selectfont Geometria dei pori e permeabilità\end{minipage}};
\node[anchor=north west,inner sep=0] at ([xshift=1.05cm,yshift=-6.4cm]current page.north west){\begin{minipage}{15.8cm}\raggedright\color{Lilac}\fontsize{15}{21}\selectfont Sezioni, carote e inferenza su un piccolo campione osservativo\end{minipage}};
\node[anchor=north west,inner sep=0] at ([xshift=1.05cm,yshift=-8.6cm]current page.north west){\begin{minipage}{16.5cm}\color{white}\fontsize{11}{15}\selectfont 21 slide · teoria, derivazioni, figure e applicazioni in R.\par\smallskip Provenienza dei dati e sorgenti nei materiali allegati.\end{minipage}};
\node[anchor=south east,inner sep=0] at ([xshift=18.95cm,yshift=-11.35cm]current page.north west){\color{Lilac!45}\bfseries\fontsize{64}{68}\selectfont P01};
\end{tikzpicture}\end{frame}
\begin{frame}[plain]\hypertarget{slide2}{}\bookmark[level=0,page=2]{Indice}\begin{tikzpicture}[remember picture,overlay]\node[anchor=north west,inner sep=0] at ([xshift=0.95cm,yshift=-0.65cm]current page.north west){\color{Navy}\bfseries\fontsize{22}{27}\selectfont Percorso del modulo};
\draw[Classical,line width=1.3pt] ([xshift=.95cm,yshift=-1.68cm]current page.north west)--([xshift=2.25cm,yshift=-1.68cm]current page.north west);\node[anchor=north west,inner sep=0] at ([xshift=0.95cm,yshift=-2.23cm]current page.north west){\color{Classical}\bfseries\fontsize{13}{16}\selectfont 3};
\node[anchor=north west,inner sep=0] at ([xshift=2.0cm,yshift=-2.23cm]current page.north west){\begin{minipage}{7.25cm}\raggedright\color{Navy}\bfseries\fontsize{12.4}{15.5}\selectfont\hyperlink{slide3}{Domanda geologica e supporto della misura}\end{minipage}};
\node[anchor=north west,inner sep=0] at ([xshift=0.95cm,yshift=-3.5300000000000002cm]current page.north west){\color{Classical}\bfseries\fontsize{13}{16}\selectfont 8};
\node[anchor=north west,inner sep=0] at ([xshift=2.0cm,yshift=-3.5300000000000002cm]current page.north west){\begin{minipage}{7.25cm}\raggedright\color{Navy}\bfseries\fontsize{12.4}{15.5}\selectfont\hyperlink{slide8}{Visualizzazione e unità di analisi}\end{minipage}};
\node[anchor=north west,inner sep=0] at ([xshift=0.95cm,yshift=-4.83cm]current page.north west){\color{Classical}\bfseries\fontsize{13}{16}\selectfont 10};
\node[anchor=north west,inner sep=0] at ([xshift=2.0cm,yshift=-4.83cm]current page.north west){\begin{minipage}{7.25cm}\raggedright\color{Navy}\bfseries\fontsize{12.4}{15.5}\selectfont\hyperlink{slide10}{Modello e interpretazione del parametro}\end{minipage}};
\node[anchor=north west,inner sep=0] at ([xshift=10.299999999999999cm,yshift=-2.23cm]current page.north west){\color{Classical}\bfseries\fontsize{13}{16}\selectfont 15};
\node[anchor=north west,inner sep=0] at ([xshift=11.35cm,yshift=-2.23cm]current page.north west){\begin{minipage}{7.25cm}\raggedright\color{Navy}\bfseries\fontsize{12.4}{15.5}\selectfont\hyperlink{slide15}{Ricampionamento e sensibilità}\end{minipage}};
\node[anchor=north west,inner sep=0] at ([xshift=10.299999999999999cm,yshift=-3.5300000000000002cm]current page.north west){\color{Classical}\bfseries\fontsize{13}{16}\selectfont 20};
\node[anchor=north west,inner sep=0] at ([xshift=11.35cm,yshift=-3.5300000000000002cm]current page.north west){\begin{minipage}{7.25cm}\raggedright\color{Navy}\bfseries\fontsize{12.4}{15.5}\selectfont\hyperlink{slide20}{Argomentazione geologica e consegna}\end{minipage}};
\node[anchor=north west,inner sep=0] at ([xshift=0.95cm,yshift=-10.55cm]current page.north west){\begin{minipage}{18cm}\color{Muted}\fontsize{10}{13}\selectfont I numeri indicano le pagine iniziali; i titoli sono cliccabili. Codice, figure e verifiche sono inclusi nel progetto.\end{minipage}};
\end{tikzpicture}\end{frame}
\begin{lrbox}{\TitleBox}\begin{minipage}[t]{18.2cm}\raggedright\sloppy\color{Ink}\fontsize{20.7}{24}\selectfont \bfseries\color{Navy}Area visibile dei pori e trasmissione del fluido\end{minipage}\end{lrbox}
\begin{lrbox}{\BodyBox}\begin{minipage}[t]{10.25cm}\raggedright\sloppy\color{Ink}\fontsize{12.4}{16.4}\selectfont La permeabilità dipende dalla struttura dei percorsi attraverso i quali un fluido attraversa la roccia, mentre una sezione bidimensionale descrive soltanto parte della geometria dei pori. Il progetto esamina l’associazione fra frazione areale dei pori e permeabilità in dodici carote di un dataset storico BP, distribuito con R. Quattro sezioni per carota forniscono quarantotto immagini, alle quali si associa la permeabilità misurata sulla carota. La domanda riguarda un’associazione osservativa entro questo campione, la cui interpretazione richiede di mantenere distinte le scale delle due misure.\end{minipage}\end{lrbox}
\begin{lrbox}{\SideBox}\begin{minipage}[t]{7.45cm}\raggedright\sloppy\color{Ink}\fontsize{12.2}{16.1}\selectfont \input{tikz/P01/D01.tex}\end{minipage}\end{lrbox}
\begin{lrbox}{\RecallBox}\begin{minipage}[t]{16.45cm}\raggedright\sloppy\color{Ink}\fontsize{10.1}{13.2}\selectfont La relazione fra un indicatore geometrico e una proprietà di trasporto richiede un’interpretazione fisica.\end{minipage}\end{lrbox}
\MeasureBox{3-TitleBox}{\TitleBox}
\MeasureBox{3-BodyBox}{\BodyBox}
\MeasureBox{3-SideBox}{\SideBox}
\MeasureBox{3-RecallBox}{\RecallBox}
\begin{frame}[plain]\hypertarget{slide3}{}
\bookmark[level=0,page=3]{Domanda geologica e supporto della misura}
\bookmark[level=1,page=3]{Area visibile dei pori e trasmissione del fluido}
\begin{tikzpicture}[remember picture,overlay]\node[anchor=north east,inner sep=0] at ([xshift=19.05cm,yshift=-0.16cm]current page.north west){\color{Classical}\fontsize{8}{10}\selectfont\hyperlink{slide2}{Indice}};
\node[anchor=north west,inner sep=0] at ([xshift=0.95cm,yshift=-0.62cm]current page.north west){\usebox{\TitleBox}};
\draw[Classical,line width=1.3pt] ([xshift=.95cm,yshift=-1.68cm]current page.north west)--([xshift=2.25cm,yshift=-1.68cm]current page.north west);\node[anchor=north west,inner sep=0] at ([xshift=0.95cm,yshift=-2.12cm]current page.north west){\usebox{\BodyBox}};
\node[anchor=north west,inner sep=0] at ([xshift=11.6cm,yshift=-2.1cm]current page.north west){\color{Quantum}\bfseries\fontsize{11.5}{14}\selectfont Lettura del diagramma};
\node[anchor=north west,inner sep=0] at ([xshift=11.6cm,yshift=-2.75cm]current page.north west){\usebox{\SideBox}};
\draw[Rule,line width=.5pt] ([xshift=.95cm,yshift=-10.43cm]current page.north west)--([xshift=19.05cm,yshift=-10.43cm]current page.north west);\node[anchor=north west,inner sep=0] at ([xshift=0.95cm,yshift=-10.66cm]current page.north west){\color{Classical}\bfseries\fontsize{9.4}{12}\selectfont Richiami};
\node[anchor=north west,inner sep=0] at ([xshift=2.6cm,yshift=-10.64cm]current page.north west){\usebox{\RecallBox}};
\draw[Rule,line width=.35pt] ([xshift=.95cm,yshift=-11.92cm]current page.north west)--([xshift=19.05cm,yshift=-11.92cm]current page.north west);\node[anchor=north west,inner sep=0] at ([xshift=0.95cm,yshift=-12.06cm]current page.north west){\color{Muted}\fontsize{7.7}{9.3}\selectfont P01 · Geometria dei pori e permeabilità};
\node[anchor=north east,inner sep=0] at ([xshift=17.65cm,yshift=-12.06cm]current page.north west){\color{Classical}\fontsize{7.3}{9.1}\selectfont R datasets: rock; BP Research, analisi immagini R. Katz};
\node[anchor=north east,inner sep=0] at ([xshift=19.05cm,yshift=-12.06cm]current page.north west){\color{Muted}\fontsize{7.7}{9.3}\selectfont 3/21};
\end{tikzpicture}\end{frame}
\begin{lrbox}{\TitleBox}\begin{minipage}[t]{18.2cm}\raggedright\sloppy\color{Ink}\fontsize{20.7}{24}\selectfont \bfseries\color{Navy}Quattro sezioni condividono una permeabilità\end{minipage}\end{lrbox}
\begin{lrbox}{\BodyBox}\begin{minipage}[t]{10.25cm}\raggedright\sloppy\color{Ink}\fontsize{12.4}{16.4}\selectfont Il manuale descrive dodici carote, ciascuna campionata attraverso quattro sezioni. La tabella distribuita contiene la permeabilità ripetuta sulle quattro righe corrispondenti alle immagini della stessa carota. Queste ripetizioni rappresentano una struttura di misura e devono essere conservate durante l’analisi. Trattare tutte le righe come risposte indipendenti attribuirebbe al campione un’informazione superiore a quella disponibile per la relazione con la permeabilità, perché la risposta non è stata misurata separatamente su ogni immagine.\end{minipage}\end{lrbox}
\begin{lrbox}{\SideBox}\begin{minipage}[t]{7.45cm}\raggedright\sloppy\color{Ink}\fontsize{12.2}{16.1}\selectfont \input{tikz/P01/D02.tex}\end{minipage}\end{lrbox}
\begin{lrbox}{\RecallBox}\begin{minipage}[t]{16.45cm}\raggedright\sloppy\color{Ink}\fontsize{10.1}{13.2}\selectfont La carota costituisce l’unità della risposta utilizzata nel modello.\end{minipage}\end{lrbox}
\MeasureBox{4-TitleBox}{\TitleBox}
\MeasureBox{4-BodyBox}{\BodyBox}
\MeasureBox{4-SideBox}{\SideBox}
\MeasureBox{4-RecallBox}{\RecallBox}
\begin{frame}[plain]\hypertarget{slide4}{}
\bookmark[level=1,page=4]{Quattro sezioni condividono una permeabilità}
\begin{tikzpicture}[remember picture,overlay]\node[anchor=north east,inner sep=0] at ([xshift=19.05cm,yshift=-0.16cm]current page.north west){\color{Classical}\fontsize{8}{10}\selectfont\hyperlink{slide2}{Indice}};
\node[anchor=north west,inner sep=0] at ([xshift=0.95cm,yshift=-0.62cm]current page.north west){\usebox{\TitleBox}};
\draw[Classical,line width=1.3pt] ([xshift=.95cm,yshift=-1.68cm]current page.north west)--([xshift=2.25cm,yshift=-1.68cm]current page.north west);\node[anchor=north west,inner sep=0] at ([xshift=0.95cm,yshift=-2.12cm]current page.north west){\usebox{\BodyBox}};
\node[anchor=north west,inner sep=0] at ([xshift=11.6cm,yshift=-2.1cm]current page.north west){\color{Quantum}\bfseries\fontsize{11.5}{14}\selectfont Lettura del diagramma};
\node[anchor=north west,inner sep=0] at ([xshift=11.6cm,yshift=-2.75cm]current page.north west){\usebox{\SideBox}};
\draw[Rule,line width=.5pt] ([xshift=.95cm,yshift=-10.43cm]current page.north west)--([xshift=19.05cm,yshift=-10.43cm]current page.north west);\node[anchor=north west,inner sep=0] at ([xshift=0.95cm,yshift=-10.66cm]current page.north west){\color{Classical}\bfseries\fontsize{9.4}{12}\selectfont Richiami};
\node[anchor=north west,inner sep=0] at ([xshift=2.6cm,yshift=-10.64cm]current page.north west){\usebox{\RecallBox}};
\draw[Rule,line width=.35pt] ([xshift=.95cm,yshift=-11.92cm]current page.north west)--([xshift=19.05cm,yshift=-11.92cm]current page.north west);\node[anchor=north west,inner sep=0] at ([xshift=0.95cm,yshift=-12.06cm]current page.north west){\color{Muted}\fontsize{7.7}{9.3}\selectfont P01 · Geometria dei pori e permeabilità};
\node[anchor=north east,inner sep=0] at ([xshift=17.65cm,yshift=-12.06cm]current page.north west){\color{Classical}\fontsize{7.3}{9.1}\selectfont };
\node[anchor=north east,inner sep=0] at ([xshift=19.05cm,yshift=-12.06cm]current page.north west){\color{Muted}\fontsize{7.7}{9.3}\selectfont 4/21};
\end{tikzpicture}\end{frame}
\begin{lrbox}{\TitleBox}\begin{minipage}[t]{18.2cm}\raggedright\sloppy\color{Ink}\fontsize{20.7}{24}\selectfont \bfseries\color{Navy}Ricostruzione degli identificativi e tracciabilità\end{minipage}\end{lrbox}
\begin{lrbox}{\BodyBox}\begin{minipage}[t]{10.25cm}\raggedright\sloppy\color{Ink}\fontsize{12.4}{16.4}\selectfont Il dataset originario non espone un identificativo di carota. La versione didattica lo ricostruisce da blocchi consecutivi di quattro righe, verificando che la permeabilità sia costante entro ogni blocco e coerente con la struttura descritta dal manuale. L’identificativo rimane una ricostruzione documentata, da confermare sui metadati originali in un’indagine applicata. Il codice confronta inoltre area e permeabilità con il dataset rock incluso in R, affinché la preparazione della tabella resti verificabile rispetto ai valori distribuiti.\end{minipage}\end{lrbox}
\begin{lrbox}{\SideBox}\begin{minipage}[t]{7.45cm}\raggedright\sloppy\color{Ink}\fontsize{12.2}{16.1}\selectfont La provenienza conserva tabella originaria, trasformazione degli identificativi e attribuzione a BP Research e Ronit Katz. Località, età e protocollo completo del campionamento restano assenti dalla tabella e limitano la definizione della popolazione alla quale estendere i risultati.\end{minipage}\end{lrbox}
\begin{lrbox}{\RecallBox}\begin{minipage}[t]{16.45cm}\raggedright\sloppy\color{Ink}\fontsize{10.1}{13.2}\selectfont Un identificativo ricostruito deve conservare la regola adottata e il limite della sua verifica.\end{minipage}\end{lrbox}
\MeasureBox{5-TitleBox}{\TitleBox}
\MeasureBox{5-BodyBox}{\BodyBox}
\MeasureBox{5-SideBox}{\SideBox}
\MeasureBox{5-RecallBox}{\RecallBox}
\begin{frame}[plain]\hypertarget{slide5}{}
\bookmark[level=1,page=5]{Ricostruzione degli identificativi e tracciabilità}
\begin{tikzpicture}[remember picture,overlay]\node[anchor=north east,inner sep=0] at ([xshift=19.05cm,yshift=-0.16cm]current page.north west){\color{Classical}\fontsize{8}{10}\selectfont\hyperlink{slide2}{Indice}};
\node[anchor=north west,inner sep=0] at ([xshift=0.95cm,yshift=-0.62cm]current page.north west){\usebox{\TitleBox}};
\draw[Classical,line width=1.3pt] ([xshift=.95cm,yshift=-1.68cm]current page.north west)--([xshift=2.25cm,yshift=-1.68cm]current page.north west);\node[anchor=north west,inner sep=0] at ([xshift=0.95cm,yshift=-2.12cm]current page.north west){\usebox{\BodyBox}};
\node[anchor=north west,inner sep=0] at ([xshift=11.6cm,yshift=-2.1cm]current page.north west){\color{Quantum}\bfseries\fontsize{11.5}{14}\selectfont Sviluppo e interpretazione};
\node[anchor=north west,inner sep=0] at ([xshift=11.6cm,yshift=-2.75cm]current page.north west){\usebox{\SideBox}};
\draw[Rule,line width=.5pt] ([xshift=.95cm,yshift=-10.43cm]current page.north west)--([xshift=19.05cm,yshift=-10.43cm]current page.north west);\node[anchor=north west,inner sep=0] at ([xshift=0.95cm,yshift=-10.66cm]current page.north west){\color{Classical}\bfseries\fontsize{9.4}{12}\selectfont Richiami};
\node[anchor=north west,inner sep=0] at ([xshift=2.6cm,yshift=-10.64cm]current page.north west){\usebox{\RecallBox}};
\draw[Rule,line width=.35pt] ([xshift=.95cm,yshift=-11.92cm]current page.north west)--([xshift=19.05cm,yshift=-11.92cm]current page.north west);\node[anchor=north west,inner sep=0] at ([xshift=0.95cm,yshift=-12.06cm]current page.north west){\color{Muted}\fontsize{7.7}{9.3}\selectfont P01 · Geometria dei pori e permeabilità};
\node[anchor=north east,inner sep=0] at ([xshift=17.65cm,yshift=-12.06cm]current page.north west){\color{Classical}\fontsize{7.3}{9.1}\selectfont };
\node[anchor=north east,inner sep=0] at ([xshift=19.05cm,yshift=-12.06cm]current page.north west){\color{Muted}\fontsize{7.7}{9.3}\selectfont 5/21};
\end{tikzpicture}\end{frame}
\begin{lrbox}{\TitleBox}\begin{minipage}[t]{18.2cm}\raggedright\sloppy\color{Ink}\fontsize{20.7}{24}\selectfont \bfseries\color{Navy}Dal conteggio dei pixel alla frazione areale\end{minipage}\end{lrbox}
\begin{lrbox}{\BodyBox}\begin{minipage}[t]{10.25cm}\raggedright\sloppy\color{Ink}\fontsize{12.4}{16.4}\selectfont L’area dei pori viene divisa per i 65536 pixel dell’immagine, pari a 256 per 256, e moltiplicata per cento. Il risultato misura una frazione areale bidimensionale secondo la segmentazione utilizzata nelle immagini. Il passaggio rende esplicito il denominatore e consente di interpretare un incremento di un punto percentuale nel modello successivo. Una conversione in porosità volumetrica o efficace richiederebbe ulteriori ipotesi geometriche e informazioni sulla rappresentatività delle sezioni e sulla connettività dei pori.\par\vspace{.28cm}{\centering\fontsize{13}{17}\selectfont\(\displaystyle\begin{gathered}A_{ij}=100\frac{a_{ij}}{256^2}\quad\text{in percentuale}\end{gathered}\)\par}\end{minipage}\end{lrbox}
\begin{lrbox}{\SideBox}\begin{minipage}[t]{7.45cm}\raggedright\sloppy\color{Ink}\fontsize{12.2}{16.1}\selectfont i = 1,…,4: sezione; j = 1,…,12: carota. $a_{ij}$ conta i pixel dei pori; $A_{ij}$ è la percentuale dell’immagine occupata.\par\vspace{.25cm}\lstinputlisting[language=R]{estratti/P01/preparazione.R}\end{minipage}\end{lrbox}
\begin{lrbox}{\RecallBox}\begin{minipage}[t]{16.45cm}\raggedright\sloppy\color{Ink}\fontsize{10.1}{13.2}\selectfont Il denominatore dell’immagine definisce la quantità geometrica che entra nell’analisi.\end{minipage}\end{lrbox}
\MeasureBox{6-TitleBox}{\TitleBox}
\MeasureBox{6-BodyBox}{\BodyBox}
\MeasureBox{6-SideBox}{\SideBox}
\MeasureBox{6-RecallBox}{\RecallBox}
\begin{frame}[plain]\hypertarget{slide6}{}
\bookmark[level=1,page=6]{Dal conteggio dei pixel alla frazione areale}
\begin{tikzpicture}[remember picture,overlay]\node[anchor=north east,inner sep=0] at ([xshift=19.05cm,yshift=-0.16cm]current page.north west){\color{Classical}\fontsize{8}{10}\selectfont\hyperlink{slide2}{Indice}};
\node[anchor=north west,inner sep=0] at ([xshift=0.95cm,yshift=-0.62cm]current page.north west){\usebox{\TitleBox}};
\draw[Classical,line width=1.3pt] ([xshift=.95cm,yshift=-1.68cm]current page.north west)--([xshift=2.25cm,yshift=-1.68cm]current page.north west);\node[anchor=north west,inner sep=0] at ([xshift=0.95cm,yshift=-2.12cm]current page.north west){\usebox{\BodyBox}};
\node[anchor=north west,inner sep=0] at ([xshift=11.6cm,yshift=-2.1cm]current page.north west){\color{Quantum}\bfseries\fontsize{11.5}{14}\selectfont Implementazione in R};
\node[anchor=north west,inner sep=0] at ([xshift=11.6cm,yshift=-2.75cm]current page.north west){\usebox{\SideBox}};
\draw[Rule,line width=.5pt] ([xshift=.95cm,yshift=-10.43cm]current page.north west)--([xshift=19.05cm,yshift=-10.43cm]current page.north west);\node[anchor=north west,inner sep=0] at ([xshift=0.95cm,yshift=-10.66cm]current page.north west){\color{Classical}\bfseries\fontsize{9.4}{12}\selectfont Richiami};
\node[anchor=north west,inner sep=0] at ([xshift=2.6cm,yshift=-10.64cm]current page.north west){\usebox{\RecallBox}};
\draw[Rule,line width=.35pt] ([xshift=.95cm,yshift=-11.92cm]current page.north west)--([xshift=19.05cm,yshift=-11.92cm]current page.north west);\node[anchor=north west,inner sep=0] at ([xshift=0.95cm,yshift=-12.06cm]current page.north west){\color{Muted}\fontsize{7.7}{9.3}\selectfont P01 · Geometria dei pori e permeabilità};
\node[anchor=north east,inner sep=0] at ([xshift=17.65cm,yshift=-12.06cm]current page.north west){\color{Classical}\fontsize{7.3}{9.1}\selectfont };
\node[anchor=north east,inner sep=0] at ([xshift=19.05cm,yshift=-12.06cm]current page.north west){\color{Muted}\fontsize{7.7}{9.3}\selectfont 6/21};
\end{tikzpicture}\end{frame}
\begin{lrbox}{\TitleBox}\begin{minipage}[t]{18.2cm}\raggedright\sloppy\color{Ink}\fontsize{20.7}{24}\selectfont \bfseries\color{Navy}Una variabile dalla definizione non verificata\end{minipage}\end{lrbox}
\begin{lrbox}{\BodyBox}\begin{minipage}[t]{10.25cm}\raggedright\sloppy\color{Ink}\fontsize{12.4}{16.4}\selectfont La colonna shape è descritta nel manuale come rapporto fra perimetro e radice quadrata dell’area, ma questa formula non riproduce i valori presenti nella tabella. Il progetto conserva la colonna originaria e ne documenta l’incongruenza, escludendola dal modello principale finché il suo significato non sia chiarito. Modificare una formula per ottenere una corrispondenza numerica potrebbe costruire una definizione priva di riscontro nel protocollo. Il controllo del dizionario precede quindi l’interpretazione di un predittore, anche quando tutti i suoi valori sono numerici.\par\vspace{.28cm}{\centering\fontsize{13}{17}\selectfont\(\displaystyle\begin{gathered}q_i^{controllo}=\frac{perimetro_i}{\sqrt{area_i}}\end{gathered}\)\par}\end{minipage}\end{lrbox}
\begin{lrbox}{\SideBox}\begin{minipage}[t]{7.45cm}\raggedright\sloppy\color{Ink}\fontsize{12.2}{16.1}\selectfont L’esercizio di approfondimento confronta la formula dichiarata con la colonna distribuita e registra lo scostamento. Il risultato del controllo giustifica una richiesta di chiarimento o l’esclusione motivata della variabile; non identifica da solo la definizione originaria.\end{minipage}\end{lrbox}
\begin{lrbox}{\RecallBox}\begin{minipage}[t]{16.45cm}\raggedright\sloppy\color{Ink}\fontsize{10.1}{13.2}\selectfont La disponibilità di una colonna numerica non garantisce che il suo significato sia ricostruibile.\end{minipage}\end{lrbox}
\MeasureBox{7-TitleBox}{\TitleBox}
\MeasureBox{7-BodyBox}{\BodyBox}
\MeasureBox{7-SideBox}{\SideBox}
\MeasureBox{7-RecallBox}{\RecallBox}
\begin{frame}[plain]\hypertarget{slide7}{}
\bookmark[level=1,page=7]{Una variabile dalla definizione non verificata}
\begin{tikzpicture}[remember picture,overlay]\node[anchor=north east,inner sep=0] at ([xshift=19.05cm,yshift=-0.16cm]current page.north west){\color{Classical}\fontsize{8}{10}\selectfont\hyperlink{slide2}{Indice}};
\node[anchor=north west,inner sep=0] at ([xshift=0.95cm,yshift=-0.62cm]current page.north west){\usebox{\TitleBox}};
\draw[Classical,line width=1.3pt] ([xshift=.95cm,yshift=-1.68cm]current page.north west)--([xshift=2.25cm,yshift=-1.68cm]current page.north west);\node[anchor=north west,inner sep=0] at ([xshift=0.95cm,yshift=-2.12cm]current page.north west){\usebox{\BodyBox}};
\node[anchor=north west,inner sep=0] at ([xshift=11.6cm,yshift=-2.1cm]current page.north west){\color{Quantum}\bfseries\fontsize{11.5}{14}\selectfont Sviluppo e interpretazione};
\node[anchor=north west,inner sep=0] at ([xshift=11.6cm,yshift=-2.75cm]current page.north west){\usebox{\SideBox}};
\draw[Rule,line width=.5pt] ([xshift=.95cm,yshift=-10.43cm]current page.north west)--([xshift=19.05cm,yshift=-10.43cm]current page.north west);\node[anchor=north west,inner sep=0] at ([xshift=0.95cm,yshift=-10.66cm]current page.north west){\color{Classical}\bfseries\fontsize{9.4}{12}\selectfont Richiami};
\node[anchor=north west,inner sep=0] at ([xshift=2.6cm,yshift=-10.64cm]current page.north west){\usebox{\RecallBox}};
\draw[Rule,line width=.35pt] ([xshift=.95cm,yshift=-11.92cm]current page.north west)--([xshift=19.05cm,yshift=-11.92cm]current page.north west);\node[anchor=north west,inner sep=0] at ([xshift=0.95cm,yshift=-12.06cm]current page.north west){\color{Muted}\fontsize{7.7}{9.3}\selectfont P01 · Geometria dei pori e permeabilità};
\node[anchor=north east,inner sep=0] at ([xshift=17.65cm,yshift=-12.06cm]current page.north west){\color{Classical}\fontsize{7.3}{9.1}\selectfont };
\node[anchor=north east,inner sep=0] at ([xshift=19.05cm,yshift=-12.06cm]current page.north west){\color{Muted}\fontsize{7.7}{9.3}\selectfont 7/21};
\end{tikzpicture}\end{frame}
\begin{lrbox}{\TitleBox}\begin{minipage}[t]{18.2cm}\raggedright\sloppy\color{Ink}\fontsize{20.7}{24}\selectfont \bfseries\color{Navy}Le immagini rivelano la ripetizione della risposta\end{minipage}\end{lrbox}
\begin{lrbox}{\BodyBox}\begin{minipage}[t]{10.25cm}\raggedright\sloppy\color{Ink}\fontsize{12.4}{16.4}\selectfont Ogni punto rappresenta una sezione, mentre la permeabilità sull’asse verticale appartiene alla carota. Le sequenze orizzontali di quattro punti rendono visibile la risposta ripetuta e la variabilità dell’area fra immagini dello stesso campione. La scala logaritmica verticale permette di leggere permeabilità distribuite su più ordini di grandezza. Il grafico serve a riconoscere la struttura delle misure; un’analisi che utilizzasse quarantotto errori indipendenti per la risposta richiederebbe un’ipotesi incompatibile con tale struttura.\end{minipage}\end{lrbox}
\begin{lrbox}{\SideBox}\begin{minipage}[t]{7.45cm}\raggedright\sloppy\color{Ink}\fontsize{12.2}{16.1}\selectfont \includegraphics[width=7.45cm]{figure/P01/F01.pdf}\par\vspace{.16cm}{\fontsize{9.5}{12.6}\selectfont Quarantotto sezioni su dodici carote. La permeabilità è in millidarcy e l’asse verticale è logaritmico.}\end{minipage}\end{lrbox}
\begin{lrbox}{\RecallBox}\begin{minipage}[t]{16.45cm}\raggedright\sloppy\color{Ink}\fontsize{10.1}{13.2}\selectfont La ripetizione della risposta nel grafico rende visibile il livello al quale è stata acquisita.\end{minipage}\end{lrbox}
\MeasureBox{8-TitleBox}{\TitleBox}
\MeasureBox{8-BodyBox}{\BodyBox}
\MeasureBox{8-SideBox}{\SideBox}
\MeasureBox{8-RecallBox}{\RecallBox}
\begin{frame}[plain]\hypertarget{slide8}{}
\bookmark[level=0,page=8]{Visualizzazione e unità di analisi}
\bookmark[level=1,page=8]{Le immagini rivelano la ripetizione della risposta}
\begin{tikzpicture}[remember picture,overlay]\node[anchor=north east,inner sep=0] at ([xshift=19.05cm,yshift=-0.16cm]current page.north west){\color{Classical}\fontsize{8}{10}\selectfont\hyperlink{slide2}{Indice}};
\node[anchor=north west,inner sep=0] at ([xshift=0.95cm,yshift=-0.62cm]current page.north west){\usebox{\TitleBox}};
\draw[Classical,line width=1.3pt] ([xshift=.95cm,yshift=-1.68cm]current page.north west)--([xshift=2.25cm,yshift=-1.68cm]current page.north west);\node[anchor=north west,inner sep=0] at ([xshift=0.95cm,yshift=-2.12cm]current page.north west){\usebox{\BodyBox}};
\node[anchor=north west,inner sep=0] at ([xshift=11.6cm,yshift=-2.1cm]current page.north west){\color{Quantum}\bfseries\fontsize{11.5}{14}\selectfont Evidenza grafica};
\node[anchor=north west,inner sep=0] at ([xshift=11.6cm,yshift=-2.75cm]current page.north west){\usebox{\SideBox}};
\draw[Rule,line width=.5pt] ([xshift=.95cm,yshift=-10.43cm]current page.north west)--([xshift=19.05cm,yshift=-10.43cm]current page.north west);\node[anchor=north west,inner sep=0] at ([xshift=0.95cm,yshift=-10.66cm]current page.north west){\color{Classical}\bfseries\fontsize{9.4}{12}\selectfont Richiami};
\node[anchor=north west,inner sep=0] at ([xshift=2.6cm,yshift=-10.64cm]current page.north west){\usebox{\RecallBox}};
\draw[Rule,line width=.35pt] ([xshift=.95cm,yshift=-11.92cm]current page.north west)--([xshift=19.05cm,yshift=-11.92cm]current page.north west);\node[anchor=north west,inner sep=0] at ([xshift=0.95cm,yshift=-12.06cm]current page.north west){\color{Muted}\fontsize{7.7}{9.3}\selectfont P01 · Geometria dei pori e permeabilità};
\node[anchor=north east,inner sep=0] at ([xshift=17.65cm,yshift=-12.06cm]current page.north west){\color{Classical}\fontsize{7.3}{9.1}\selectfont };
\node[anchor=north east,inner sep=0] at ([xshift=19.05cm,yshift=-12.06cm]current page.north west){\color{Muted}\fontsize{7.7}{9.3}\selectfont 8/21};
\end{tikzpicture}\end{frame}
\begin{lrbox}{\TitleBox}\begin{minipage}[t]{18.2cm}\raggedright\sloppy\color{Ink}\fontsize{20.7}{24}\selectfont \bfseries\color{Navy}Una sintesi areale per ogni carota\end{minipage}\end{lrbox}
\begin{lrbox}{\BodyBox}\begin{minipage}[t]{10.25cm}\raggedright\sloppy\color{Ink}\fontsize{12.4}{16.4}\selectfont La media delle quattro frazioni areali rappresenta ciascuna carota attraverso un unico predittore, che viene associato alla sua permeabilità. Si ottengono dodici coppie, equiponderate nel modello successivo. La sintesi riduce la ripetizione della risposta, pur lasciando aperte l’incertezza della misura areale e la rappresentatività delle sezioni. I numeri accanto ai punti conservano gli identificativi ricostruiti e permettono di seguire i campioni nelle analisi di influenza, collegando il grafico alla tabella delle misure.\par\vspace{.28cm}{\centering\fontsize{13}{17}\selectfont\(\displaystyle\begin{gathered}\overline A_j=\frac14\sum_{i=1}^{4}A_{ij},\qquad j=1,\ldots,12\end{gathered}\)\par}\end{minipage}\end{lrbox}
\begin{lrbox}{\SideBox}\begin{minipage}[t]{7.45cm}\raggedright\sloppy\color{Ink}\fontsize{12.2}{16.1}\selectfont \includegraphics[width=7.45cm]{figure/P01/F02.pdf}\par\vspace{.16cm}{\fontsize{9.5}{12.6}\selectfont $A_{ij}$: percentuale areale della sezione i nella carota j; $\bar A_j$: media delle quattro sezioni. Ordinata: $\log_{10}(K/1\,\mathrm{mD})$.}\end{minipage}\end{lrbox}
\begin{lrbox}{\RecallBox}\begin{minipage}[t]{16.45cm}\raggedright\sloppy\color{Ink}\fontsize{10.1}{13.2}\selectfont La sintesi esplicita l’unità del modello e conserva la provenienza dei valori aggregati.\end{minipage}\end{lrbox}
\MeasureBox{9-TitleBox}{\TitleBox}
\MeasureBox{9-BodyBox}{\BodyBox}
\MeasureBox{9-SideBox}{\SideBox}
\MeasureBox{9-RecallBox}{\RecallBox}
\begin{frame}[plain]\hypertarget{slide9}{}
\bookmark[level=1,page=9]{Una sintesi areale per ogni carota}
\begin{tikzpicture}[remember picture,overlay]\node[anchor=north east,inner sep=0] at ([xshift=19.05cm,yshift=-0.16cm]current page.north west){\color{Classical}\fontsize{8}{10}\selectfont\hyperlink{slide2}{Indice}};
\node[anchor=north west,inner sep=0] at ([xshift=0.95cm,yshift=-0.62cm]current page.north west){\usebox{\TitleBox}};
\draw[Classical,line width=1.3pt] ([xshift=.95cm,yshift=-1.68cm]current page.north west)--([xshift=2.25cm,yshift=-1.68cm]current page.north west);\node[anchor=north west,inner sep=0] at ([xshift=0.95cm,yshift=-2.12cm]current page.north west){\usebox{\BodyBox}};
\node[anchor=north west,inner sep=0] at ([xshift=11.6cm,yshift=-2.1cm]current page.north west){\color{Quantum}\bfseries\fontsize{11.5}{14}\selectfont Evidenza grafica};
\node[anchor=north west,inner sep=0] at ([xshift=11.6cm,yshift=-2.75cm]current page.north west){\usebox{\SideBox}};
\draw[Rule,line width=.5pt] ([xshift=.95cm,yshift=-10.43cm]current page.north west)--([xshift=19.05cm,yshift=-10.43cm]current page.north west);\node[anchor=north west,inner sep=0] at ([xshift=0.95cm,yshift=-10.66cm]current page.north west){\color{Classical}\bfseries\fontsize{9.4}{12}\selectfont Richiami};
\node[anchor=north west,inner sep=0] at ([xshift=2.6cm,yshift=-10.64cm]current page.north west){\usebox{\RecallBox}};
\draw[Rule,line width=.35pt] ([xshift=.95cm,yshift=-11.92cm]current page.north west)--([xshift=19.05cm,yshift=-11.92cm]current page.north west);\node[anchor=north west,inner sep=0] at ([xshift=0.95cm,yshift=-12.06cm]current page.north west){\color{Muted}\fontsize{7.7}{9.3}\selectfont P01 · Geometria dei pori e permeabilità};
\node[anchor=north east,inner sep=0] at ([xshift=17.65cm,yshift=-12.06cm]current page.north west){\color{Classical}\fontsize{7.3}{9.1}\selectfont };
\node[anchor=north east,inner sep=0] at ([xshift=19.05cm,yshift=-12.06cm]current page.north west){\color{Muted}\fontsize{7.7}{9.3}\selectfont 9/21};
\end{tikzpicture}\end{frame}
\begin{lrbox}{\TitleBox}\begin{minipage}[t]{18.2cm}\raggedright\sloppy\color{Ink}\fontsize{20.7}{24}\selectfont \bfseries\color{Navy}Il modello definisce l’effetto da stimare\end{minipage}\end{lrbox}
\begin{lrbox}{\BodyBox}\begin{minipage}[t]{10.25cm}\raggedright\sloppy\color{Ink}\fontsize{12.4}{16.4}\selectfont La regressione collega il logaritmo decimale della permeabilità alla frazione areale media. La pendenza misura la variazione della risposta logaritmica per un punto percentuale di area, mentre la sua trasformazione in potenza di dieci esprime un fattore fra medie geometriche condizionate. La base del logaritmo e l’unità di riferimento devono essere dichiarate, poiché entrano nell’interpretazione dei coefficienti.\par\vspace{.28cm}{\centering\fontsize{13}{17}\selectfont\(\displaystyle\begin{gathered}Y_j=\log_{10}\!\left(\frac{K_j}{1\ \mathrm{mD}}\right)=\alpha+\beta\overline A_j+\varepsilon_j\\\frac{G(K\mid A+1)}{G(K\mid A)}=10^\beta\end{gathered}\)\par}\end{minipage}\end{lrbox}
\begin{lrbox}{\SideBox}\begin{minipage}[t]{7.45cm}\raggedright\sloppy\color{Ink}\fontsize{12.2}{16.1}\selectfont $K_j$: permeabilità della carota j; $\bar A_j$: percentuale areale media; $\alpha$: intercetta; $\beta$: pendenza. Si assume $E(\varepsilon_j\mid\bar A_j)=0$. G indica la media geometrica condizionata, definita tramite la media del logaritmo. La trasformazione inversa della retta non coincide in generale con la media aritmetica. L’associazione osservativa non identifica un effetto causale.\end{minipage}\end{lrbox}
\begin{lrbox}{\RecallBox}\begin{minipage}[t]{16.45cm}\raggedright\sloppy\color{Ink}\fontsize{10.1}{13.2}\selectfont La scala della risposta determina il significato fisico della pendenza e della trasformazione inversa.\end{minipage}\end{lrbox}
\MeasureBox{10-TitleBox}{\TitleBox}
\MeasureBox{10-BodyBox}{\BodyBox}
\MeasureBox{10-SideBox}{\SideBox}
\MeasureBox{10-RecallBox}{\RecallBox}
\begin{frame}[plain]\hypertarget{slide10}{}
\bookmark[level=0,page=10]{Modello e interpretazione del parametro}
\bookmark[level=1,page=10]{Il modello definisce l’effetto da stimare}
\begin{tikzpicture}[remember picture,overlay]\node[anchor=north east,inner sep=0] at ([xshift=19.05cm,yshift=-0.16cm]current page.north west){\color{Classical}\fontsize{8}{10}\selectfont\hyperlink{slide2}{Indice}};
\node[anchor=north west,inner sep=0] at ([xshift=0.95cm,yshift=-0.62cm]current page.north west){\usebox{\TitleBox}};
\draw[Classical,line width=1.3pt] ([xshift=.95cm,yshift=-1.68cm]current page.north west)--([xshift=2.25cm,yshift=-1.68cm]current page.north west);\node[anchor=north west,inner sep=0] at ([xshift=0.95cm,yshift=-2.12cm]current page.north west){\usebox{\BodyBox}};
\node[anchor=north west,inner sep=0] at ([xshift=11.6cm,yshift=-2.1cm]current page.north west){\color{Quantum}\bfseries\fontsize{11.5}{14}\selectfont Sviluppo e interpretazione};
\node[anchor=north west,inner sep=0] at ([xshift=11.6cm,yshift=-2.75cm]current page.north west){\usebox{\SideBox}};
\draw[Rule,line width=.5pt] ([xshift=.95cm,yshift=-10.43cm]current page.north west)--([xshift=19.05cm,yshift=-10.43cm]current page.north west);\node[anchor=north west,inner sep=0] at ([xshift=0.95cm,yshift=-10.66cm]current page.north west){\color{Classical}\bfseries\fontsize{9.4}{12}\selectfont Richiami};
\node[anchor=north west,inner sep=0] at ([xshift=2.6cm,yshift=-10.64cm]current page.north west){\usebox{\RecallBox}};
\draw[Rule,line width=.35pt] ([xshift=.95cm,yshift=-11.92cm]current page.north west)--([xshift=19.05cm,yshift=-11.92cm]current page.north west);\node[anchor=north west,inner sep=0] at ([xshift=0.95cm,yshift=-12.06cm]current page.north west){\color{Muted}\fontsize{7.7}{9.3}\selectfont P01 · Geometria dei pori e permeabilità};
\node[anchor=north east,inner sep=0] at ([xshift=17.65cm,yshift=-12.06cm]current page.north west){\color{Classical}\fontsize{7.3}{9.1}\selectfont };
\node[anchor=north east,inner sep=0] at ([xshift=19.05cm,yshift=-12.06cm]current page.north west){\color{Muted}\fontsize{7.7}{9.3}\selectfont 10/21};
\end{tikzpicture}\end{frame}
\begin{lrbox}{\TitleBox}\begin{minipage}[t]{18.2cm}\raggedright\sloppy\color{Ink}\fontsize{20.7}{24}\selectfont \bfseries\color{Navy}Stima della pendenza attraverso la covariazione\end{minipage}\end{lrbox}
\begin{lrbox}{\BodyBox}\begin{minipage}[t]{10.25cm}\raggedright\sloppy\color{Ink}\fontsize{12.4}{16.4}\selectfont La stima ai minimi quadrati confronta la covariazione fra area e risposta logaritmica con la dispersione dell’area. Centrare entrambe le variabili rende esplicito che la pendenza descrive come i loro scostamenti si associno nel campione. Il codice verifica l’equivalenza fra la formula e il coefficiente restituito da lm. La stima ottenuta è circa meno 0,08554 per punto percentuale, mentre il fattore in scala originaria è circa 0,8212. Questi valori richiedono un intervallo e una diagnosi del modello prima dell’interpretazione conclusiva.\par\vspace{.28cm}{\centering\fontsize{13}{17}\selectfont\(\displaystyle\begin{gathered}\widehat\beta=\frac{\sum_j(\overline A_j-\overline A)(Y_j-\overline Y)}{\sum_j(\overline A_j-\overline A)^2}\end{gathered}\)\par}\end{minipage}\end{lrbox}
\begin{lrbox}{\SideBox}\begin{minipage}[t]{7.45cm}\raggedright\sloppy\color{Ink}\fontsize{12.2}{16.1}\selectfont \input{tikz/P01/D03.tex}\par\vspace{.16cm}{\fontsize{9.5}{12.6}\selectfont $\bar A_j$: area media della carota j; $Y_j=\log_{10}(K_j/1\,\mathrm{mD})$. $\bar A$ e $\bar Y$ sono le medie sulle dodici carote.}\end{minipage}\end{lrbox}
\begin{lrbox}{\RecallBox}\begin{minipage}[t]{16.45cm}\raggedright\sloppy\color{Ink}\fontsize{10.1}{13.2}\selectfont Stima e interpretazione devono conservare scala, unità e dominio dell’associazione.\end{minipage}\end{lrbox}
\MeasureBox{11-TitleBox}{\TitleBox}
\MeasureBox{11-BodyBox}{\BodyBox}
\MeasureBox{11-SideBox}{\SideBox}
\MeasureBox{11-RecallBox}{\RecallBox}
\begin{frame}[plain]\hypertarget{slide11}{}
\bookmark[level=1,page=11]{Stima della pendenza attraverso la covariazione}
\begin{tikzpicture}[remember picture,overlay]\node[anchor=north east,inner sep=0] at ([xshift=19.05cm,yshift=-0.16cm]current page.north west){\color{Classical}\fontsize{8}{10}\selectfont\hyperlink{slide2}{Indice}};
\node[anchor=north west,inner sep=0] at ([xshift=0.95cm,yshift=-0.62cm]current page.north west){\usebox{\TitleBox}};
\draw[Classical,line width=1.3pt] ([xshift=.95cm,yshift=-1.68cm]current page.north west)--([xshift=2.25cm,yshift=-1.68cm]current page.north west);\node[anchor=north west,inner sep=0] at ([xshift=0.95cm,yshift=-2.12cm]current page.north west){\usebox{\BodyBox}};
\node[anchor=north west,inner sep=0] at ([xshift=11.6cm,yshift=-2.1cm]current page.north west){\color{Quantum}\bfseries\fontsize{11.5}{14}\selectfont Lettura del diagramma};
\node[anchor=north west,inner sep=0] at ([xshift=11.6cm,yshift=-2.75cm]current page.north west){\usebox{\SideBox}};
\draw[Rule,line width=.5pt] ([xshift=.95cm,yshift=-10.43cm]current page.north west)--([xshift=19.05cm,yshift=-10.43cm]current page.north west);\node[anchor=north west,inner sep=0] at ([xshift=0.95cm,yshift=-10.66cm]current page.north west){\color{Classical}\bfseries\fontsize{9.4}{12}\selectfont Richiami};
\node[anchor=north west,inner sep=0] at ([xshift=2.6cm,yshift=-10.64cm]current page.north west){\usebox{\RecallBox}};
\draw[Rule,line width=.35pt] ([xshift=.95cm,yshift=-11.92cm]current page.north west)--([xshift=19.05cm,yshift=-11.92cm]current page.north west);\node[anchor=north west,inner sep=0] at ([xshift=0.95cm,yshift=-12.06cm]current page.north west){\color{Muted}\fontsize{7.7}{9.3}\selectfont P01 · Geometria dei pori e permeabilità};
\node[anchor=north east,inner sep=0] at ([xshift=17.65cm,yshift=-12.06cm]current page.north west){\color{Classical}\fontsize{7.3}{9.1}\selectfont };
\node[anchor=north east,inner sep=0] at ([xshift=19.05cm,yshift=-12.06cm]current page.north west){\color{Muted}\fontsize{7.7}{9.3}\selectfont 11/21};
\end{tikzpicture}\end{frame}
\begin{lrbox}{\TitleBox}\begin{minipage}[t]{18.2cm}\raggedright\sloppy\color{Ink}\fontsize{20.7}{24}\selectfont \bfseries\color{Navy}Adattamento e controllo numerico in R\end{minipage}\end{lrbox}
\begin{lrbox}{\BodyBox}\begin{minipage}[t]{10.25cm}\raggedright\sloppy\color{Ink}\fontsize{12.4}{16.4}\selectfont Il frammento ricostruisce la sintesi per carota, trasforma la risposta e adatta la regressione. Il confronto con la formula della covarianza costituisce un controllo indipendente dell’implementazione del coefficiente. L’uscita di summary distingue stima, errore standard, statistica e probabilità di riferimento, mentre confint restituisce l’intervallo classico sotto le ipotesi del modello. La lettura dell’output richiede di collegare queste quantità alla struttura del campione, anziché considerare il completamento del comando una verifica delle assunzioni.\end{minipage}\end{lrbox}
\begin{lrbox}{\SideBox}\begin{minipage}[t]{7.45cm}\raggedright\sloppy\color{Ink}\fontsize{12.2}{16.1}\selectfont \par\vspace{.25cm}\lstinputlisting[language=R]{estratti/P01/modello.R}\end{minipage}\end{lrbox}
\begin{lrbox}{\RecallBox}\begin{minipage}[t]{16.45cm}\raggedright\sloppy\color{Ink}\fontsize{10.1}{13.2}\selectfont Il controllo numerico verifica il calcolo del modello; le sue assunzioni richiedono evidenza ulteriore.\end{minipage}\end{lrbox}
\MeasureBox{12-TitleBox}{\TitleBox}
\MeasureBox{12-BodyBox}{\BodyBox}
\MeasureBox{12-SideBox}{\SideBox}
\MeasureBox{12-RecallBox}{\RecallBox}
\begin{frame}[plain]\hypertarget{slide12}{}
\bookmark[level=1,page=12]{Adattamento e controllo numerico in R}
\begin{tikzpicture}[remember picture,overlay]\node[anchor=north east,inner sep=0] at ([xshift=19.05cm,yshift=-0.16cm]current page.north west){\color{Classical}\fontsize{8}{10}\selectfont\hyperlink{slide2}{Indice}};
\node[anchor=north west,inner sep=0] at ([xshift=0.95cm,yshift=-0.62cm]current page.north west){\usebox{\TitleBox}};
\draw[Classical,line width=1.3pt] ([xshift=.95cm,yshift=-1.68cm]current page.north west)--([xshift=2.25cm,yshift=-1.68cm]current page.north west);\node[anchor=north west,inner sep=0] at ([xshift=0.95cm,yshift=-2.12cm]current page.north west){\usebox{\BodyBox}};
\node[anchor=north west,inner sep=0] at ([xshift=11.6cm,yshift=-2.1cm]current page.north west){\color{Quantum}\bfseries\fontsize{11.5}{14}\selectfont Implementazione in R};
\node[anchor=north west,inner sep=0] at ([xshift=11.6cm,yshift=-2.75cm]current page.north west){\usebox{\SideBox}};
\draw[Rule,line width=.5pt] ([xshift=.95cm,yshift=-10.43cm]current page.north west)--([xshift=19.05cm,yshift=-10.43cm]current page.north west);\node[anchor=north west,inner sep=0] at ([xshift=0.95cm,yshift=-10.66cm]current page.north west){\color{Classical}\bfseries\fontsize{9.4}{12}\selectfont Richiami};
\node[anchor=north west,inner sep=0] at ([xshift=2.6cm,yshift=-10.64cm]current page.north west){\usebox{\RecallBox}};
\draw[Rule,line width=.35pt] ([xshift=.95cm,yshift=-11.92cm]current page.north west)--([xshift=19.05cm,yshift=-11.92cm]current page.north west);\node[anchor=north west,inner sep=0] at ([xshift=0.95cm,yshift=-12.06cm]current page.north west){\color{Muted}\fontsize{7.7}{9.3}\selectfont P01 · Geometria dei pori e permeabilità};
\node[anchor=north east,inner sep=0] at ([xshift=17.65cm,yshift=-12.06cm]current page.north west){\color{Classical}\fontsize{7.3}{9.1}\selectfont };
\node[anchor=north east,inner sep=0] at ([xshift=19.05cm,yshift=-12.06cm]current page.north west){\color{Muted}\fontsize{7.7}{9.3}\selectfont 12/21};
\end{tikzpicture}\end{frame}
\begin{lrbox}{\TitleBox}\begin{minipage}[t]{18.2cm}\raggedright\sloppy\color{Ink}\fontsize{20.7}{24}\selectfont \bfseries\color{Navy}Intervallo classico e gradi di libertà\end{minipage}\end{lrbox}
\begin{lrbox}{\BodyBox}\begin{minipage}[t]{10.25cm}\raggedright\sloppy\color{Ink}\fontsize{12.4}{16.4}\selectfont Con dodici carote e due coefficienti stimati, la varianza residua utilizza dieci gradi di libertà. L’intervallo t per la pendenza assume errori indipendenti, gaussiani e con varianza costante attorno a una relazione lineare correttamente specificata. Nel campione l’errore standard è circa 0,06288 e l’intervallo al 95\% si estende da meno 0,22565 a più 0,05458. La presenza dello zero segnala la compatibilità di pendenze differenti con queste assunzioni e con l’informazione osservata.\par\vspace{.28cm}{\centering\fontsize{13}{17}\selectfont\(\displaystyle\begin{gathered}\widehat\sigma^2=\frac{\sum_j\widehat\varepsilon_j^2}{12-2},\qquad SE(\widehat\beta)=\sqrt{\frac{\widehat\sigma^2}{\sum_j(\overline A_j-\overline A)^2}}\\IC_{95\%}=\widehat\beta\pm t_{0{,}975;10}SE(\widehat\beta)\end{gathered}\)\par}\end{minipage}\end{lrbox}
\begin{lrbox}{\SideBox}\begin{minipage}[t]{7.45cm}\raggedright\sloppy\color{Ink}\fontsize{12.2}{16.1}\selectfont $\hat\varepsilon_j=Y_j-\hat Y_j$: residuo logaritmico; $\bar A_j$: percentuale areale della carota; $\bar A$: media delle dodici carote. SE indica l’errore standard; $t_{0,975;10}\simeq2,228$ è il quantile t. L’esattezza in piccolo campione dipende dalle ipotesi sugli errori; selezione delle carote e incertezza del predittore richiedono valutazioni ulteriori.\end{minipage}\end{lrbox}
\begin{lrbox}{\RecallBox}\begin{minipage}[t]{16.45cm}\raggedright\sloppy\color{Ink}\fontsize{10.1}{13.2}\selectfont Un intervallo acquista significato attraverso il modello probabilistico che ne giustifica la costruzione.\end{minipage}\end{lrbox}
\MeasureBox{13-TitleBox}{\TitleBox}
\MeasureBox{13-BodyBox}{\BodyBox}
\MeasureBox{13-SideBox}{\SideBox}
\MeasureBox{13-RecallBox}{\RecallBox}
\begin{frame}[plain]\hypertarget{slide13}{}
\bookmark[level=1,page=13]{Intervallo classico e gradi di libertà}
\begin{tikzpicture}[remember picture,overlay]\node[anchor=north east,inner sep=0] at ([xshift=19.05cm,yshift=-0.16cm]current page.north west){\color{Classical}\fontsize{8}{10}\selectfont\hyperlink{slide2}{Indice}};
\node[anchor=north west,inner sep=0] at ([xshift=0.95cm,yshift=-0.62cm]current page.north west){\usebox{\TitleBox}};
\draw[Classical,line width=1.3pt] ([xshift=.95cm,yshift=-1.68cm]current page.north west)--([xshift=2.25cm,yshift=-1.68cm]current page.north west);\node[anchor=north west,inner sep=0] at ([xshift=0.95cm,yshift=-2.12cm]current page.north west){\usebox{\BodyBox}};
\node[anchor=north west,inner sep=0] at ([xshift=11.6cm,yshift=-2.1cm]current page.north west){\color{Quantum}\bfseries\fontsize{11.5}{14}\selectfont Sviluppo e interpretazione};
\node[anchor=north west,inner sep=0] at ([xshift=11.6cm,yshift=-2.75cm]current page.north west){\usebox{\SideBox}};
\draw[Rule,line width=.5pt] ([xshift=.95cm,yshift=-10.43cm]current page.north west)--([xshift=19.05cm,yshift=-10.43cm]current page.north west);\node[anchor=north west,inner sep=0] at ([xshift=0.95cm,yshift=-10.66cm]current page.north west){\color{Classical}\bfseries\fontsize{9.4}{12}\selectfont Richiami};
\node[anchor=north west,inner sep=0] at ([xshift=2.6cm,yshift=-10.64cm]current page.north west){\usebox{\RecallBox}};
\draw[Rule,line width=.35pt] ([xshift=.95cm,yshift=-11.92cm]current page.north west)--([xshift=19.05cm,yshift=-11.92cm]current page.north west);\node[anchor=north west,inner sep=0] at ([xshift=0.95cm,yshift=-12.06cm]current page.north west){\color{Muted}\fontsize{7.7}{9.3}\selectfont P01 · Geometria dei pori e permeabilità};
\node[anchor=north east,inner sep=0] at ([xshift=17.65cm,yshift=-12.06cm]current page.north west){\color{Classical}\fontsize{7.3}{9.1}\selectfont };
\node[anchor=north east,inner sep=0] at ([xshift=19.05cm,yshift=-12.06cm]current page.north west){\color{Muted}\fontsize{7.7}{9.3}\selectfont 13/21};
\end{tikzpicture}\end{frame}
\begin{lrbox}{\TitleBox}\begin{minipage}[t]{18.2cm}\raggedright\sloppy\color{Ink}\fontsize{20.7}{24}\selectfont \bfseries\color{Navy}Retta e incertezza della risposta media\end{minipage}\end{lrbox}
\begin{lrbox}{\BodyBox}\begin{minipage}[t]{10.25cm}\raggedright\sloppy\color{Ink}\fontsize{12.4}{16.4}\selectfont La fascia attorno alla retta rappresenta l’intervallo puntuale al 95\% della media della risposta logaritmica, calcolato secondo il modello classico. La larghezza aumenta allontanandosi dall’area media del campione, perché diminuisce l’informazione disponibile sulla posizione della retta. Una previsione per una nuova carota richiederebbe anche la variabilità residua individuale e produrrebbe un intervallo più ampio. La figura mantiene il dominio delle aree osservate, evitando di attribuire alla relazione una validità esterna non verificata.\end{minipage}\end{lrbox}
\begin{lrbox}{\SideBox}\begin{minipage}[t]{7.45cm}\raggedright\sloppy\color{Ink}\fontsize{12.2}{16.1}\selectfont \includegraphics[width=7.45cm]{figure/P01/F03.pdf}\par\vspace{.16cm}{\fontsize{9.5}{12.6}\selectfont Retta OLS e intervallo puntuale della media in scala log10. I dodici punti conservano il livello della carota.}\end{minipage}\end{lrbox}
\begin{lrbox}{\RecallBox}\begin{minipage}[t]{16.45cm}\raggedright\sloppy\color{Ink}\fontsize{10.1}{13.2}\selectfont L’incertezza della media stimata e quella di una nuova osservazione corrispondono a quantità diverse.\end{minipage}\end{lrbox}
\MeasureBox{14-TitleBox}{\TitleBox}
\MeasureBox{14-BodyBox}{\BodyBox}
\MeasureBox{14-SideBox}{\SideBox}
\MeasureBox{14-RecallBox}{\RecallBox}
\begin{frame}[plain]\hypertarget{slide14}{}
\bookmark[level=1,page=14]{Retta e incertezza della risposta media}
\begin{tikzpicture}[remember picture,overlay]\node[anchor=north east,inner sep=0] at ([xshift=19.05cm,yshift=-0.16cm]current page.north west){\color{Classical}\fontsize{8}{10}\selectfont\hyperlink{slide2}{Indice}};
\node[anchor=north west,inner sep=0] at ([xshift=0.95cm,yshift=-0.62cm]current page.north west){\usebox{\TitleBox}};
\draw[Classical,line width=1.3pt] ([xshift=.95cm,yshift=-1.68cm]current page.north west)--([xshift=2.25cm,yshift=-1.68cm]current page.north west);\node[anchor=north west,inner sep=0] at ([xshift=0.95cm,yshift=-2.12cm]current page.north west){\usebox{\BodyBox}};
\node[anchor=north west,inner sep=0] at ([xshift=11.6cm,yshift=-2.1cm]current page.north west){\color{Quantum}\bfseries\fontsize{11.5}{14}\selectfont Evidenza grafica};
\node[anchor=north west,inner sep=0] at ([xshift=11.6cm,yshift=-2.75cm]current page.north west){\usebox{\SideBox}};
\draw[Rule,line width=.5pt] ([xshift=.95cm,yshift=-10.43cm]current page.north west)--([xshift=19.05cm,yshift=-10.43cm]current page.north west);\node[anchor=north west,inner sep=0] at ([xshift=0.95cm,yshift=-10.66cm]current page.north west){\color{Classical}\bfseries\fontsize{9.4}{12}\selectfont Richiami};
\node[anchor=north west,inner sep=0] at ([xshift=2.6cm,yshift=-10.64cm]current page.north west){\usebox{\RecallBox}};
\draw[Rule,line width=.35pt] ([xshift=.95cm,yshift=-11.92cm]current page.north west)--([xshift=19.05cm,yshift=-11.92cm]current page.north west);\node[anchor=north west,inner sep=0] at ([xshift=0.95cm,yshift=-12.06cm]current page.north west){\color{Muted}\fontsize{7.7}{9.3}\selectfont P01 · Geometria dei pori e permeabilità};
\node[anchor=north east,inner sep=0] at ([xshift=17.65cm,yshift=-12.06cm]current page.north west){\color{Classical}\fontsize{7.3}{9.1}\selectfont };
\node[anchor=north east,inner sep=0] at ([xshift=19.05cm,yshift=-12.06cm]current page.north west){\color{Muted}\fontsize{7.7}{9.3}\selectfont 14/21};
\end{tikzpicture}\end{frame}
\begin{lrbox}{\TitleBox}\begin{minipage}[t]{18.2cm}\raggedright\sloppy\color{Ink}\fontsize{20.7}{24}\selectfont \bfseries\color{Navy}Il bootstrap ricampiona intere carote\end{minipage}\end{lrbox}
\begin{lrbox}{\BodyBox}\begin{minipage}[t]{10.25cm}\raggedright\sloppy\color{Ink}\fontsize{12.4}{16.4}\selectfont Ogni replica estrae dodici carote con reinserimento e ristima la pendenza sulle coppie area–permeabilità selezionate. Le quattro sezioni sono già rappresentate dalla media della rispettiva carota, per cui il ricampionamento mantiene l’unità della risposta. Il progetto conserva quattromila piani di estrazione condivisi, che rendono confrontabili le implementazioni indipendentemente dai generatori casuali. Con dodici unità osservate, la distribuzione bootstrap esplora la sensibilità alle molteplicità del campione disponibile e richiede cautela nelle conclusioni sulla coda.\end{minipage}\end{lrbox}
\begin{lrbox}{\SideBox}\begin{minipage}[t]{7.45cm}\raggedright\sloppy\color{Ink}\fontsize{12.2}{16.1}\selectfont \input{tikz/P01/D04.tex}\end{minipage}\end{lrbox}
\begin{lrbox}{\RecallBox}\begin{minipage}[t]{16.45cm}\raggedright\sloppy\color{Ink}\fontsize{10.1}{13.2}\selectfont La scelta dell’unità ricampionata conserva la struttura sulla quale si fonda l’inferenza.\end{minipage}\end{lrbox}
\MeasureBox{15-TitleBox}{\TitleBox}
\MeasureBox{15-BodyBox}{\BodyBox}
\MeasureBox{15-SideBox}{\SideBox}
\MeasureBox{15-RecallBox}{\RecallBox}
\begin{frame}[plain]\hypertarget{slide15}{}
\bookmark[level=0,page=15]{Ricampionamento e sensibilità}
\bookmark[level=1,page=15]{Il bootstrap ricampiona intere carote}
\begin{tikzpicture}[remember picture,overlay]\node[anchor=north east,inner sep=0] at ([xshift=19.05cm,yshift=-0.16cm]current page.north west){\color{Classical}\fontsize{8}{10}\selectfont\hyperlink{slide2}{Indice}};
\node[anchor=north west,inner sep=0] at ([xshift=0.95cm,yshift=-0.62cm]current page.north west){\usebox{\TitleBox}};
\draw[Classical,line width=1.3pt] ([xshift=.95cm,yshift=-1.68cm]current page.north west)--([xshift=2.25cm,yshift=-1.68cm]current page.north west);\node[anchor=north west,inner sep=0] at ([xshift=0.95cm,yshift=-2.12cm]current page.north west){\usebox{\BodyBox}};
\node[anchor=north west,inner sep=0] at ([xshift=11.6cm,yshift=-2.1cm]current page.north west){\color{Quantum}\bfseries\fontsize{11.5}{14}\selectfont Lettura del diagramma};
\node[anchor=north west,inner sep=0] at ([xshift=11.6cm,yshift=-2.75cm]current page.north west){\usebox{\SideBox}};
\draw[Rule,line width=.5pt] ([xshift=.95cm,yshift=-10.43cm]current page.north west)--([xshift=19.05cm,yshift=-10.43cm]current page.north west);\node[anchor=north west,inner sep=0] at ([xshift=0.95cm,yshift=-10.66cm]current page.north west){\color{Classical}\bfseries\fontsize{9.4}{12}\selectfont Richiami};
\node[anchor=north west,inner sep=0] at ([xshift=2.6cm,yshift=-10.64cm]current page.north west){\usebox{\RecallBox}};
\draw[Rule,line width=.35pt] ([xshift=.95cm,yshift=-11.92cm]current page.north west)--([xshift=19.05cm,yshift=-11.92cm]current page.north west);\node[anchor=north west,inner sep=0] at ([xshift=0.95cm,yshift=-12.06cm]current page.north west){\color{Muted}\fontsize{7.7}{9.3}\selectfont P01 · Geometria dei pori e permeabilità};
\node[anchor=north east,inner sep=0] at ([xshift=17.65cm,yshift=-12.06cm]current page.north west){\color{Classical}\fontsize{7.3}{9.1}\selectfont };
\node[anchor=north east,inner sep=0] at ([xshift=19.05cm,yshift=-12.06cm]current page.north west){\color{Muted}\fontsize{7.7}{9.3}\selectfont 15/21};
\end{tikzpicture}\end{frame}
\begin{lrbox}{\TitleBox}\begin{minipage}[t]{18.2cm}\raggedright\sloppy\color{Ink}\fontsize{20.7}{24}\selectfont \bfseries\color{Navy}Distribuzione bootstrap della pendenza\end{minipage}\end{lrbox}
\begin{lrbox}{\BodyBox}\begin{minipage}[t]{10.25cm}\raggedright\sloppy\color{Ink}\fontsize{12.4}{16.4}\selectfont Le quattromila stime mostrano la variabilità prodotta dalle diverse molteplicità delle carote estratte. L’intervallo percentile utilizza i quantili al 2,5\% e 97,5\%, che in questa esecuzione valgono circa meno 0,20194 e meno 0,00893. Entrambi gli estremi risultano negativi, diversamente dall’intervallo t. La distribuzione riflette il campione disponibile e la procedura di ricampionamento; il numero elevato di repliche riduce l’errore Monte Carlo, senza aumentare il numero di carote osservate.\par\vspace{.28cm}{\centering\fontsize{13}{17}\selectfont\(\displaystyle\begin{gathered}IC_{perc}=\left[Q_{0{,}025}(\widehat\beta^*),Q_{0{,}975}(\widehat\beta^*)\right]\end{gathered}\)\par}\end{minipage}\end{lrbox}
\begin{lrbox}{\SideBox}\begin{minipage}[t]{7.45cm}\raggedright\sloppy\color{Ink}\fontsize{12.2}{16.1}\selectfont \includegraphics[width=7.45cm]{figure/P01/F04.pdf}\par\vspace{.16cm}{\fontsize{9.5}{12.6}\selectfont $\hat\beta^*$: pendenza ricampionata; $Q_p$: quantile empirico. 4000 repliche; viola: estremi percentili; ocra: zero.}\end{minipage}\end{lrbox}
\begin{lrbox}{\RecallBox}\begin{minipage}[t]{16.45cm}\raggedright\sloppy\color{Ink}\fontsize{10.1}{13.2}\selectfont Il numero di ricampionamenti controlla la simulazione della procedura e lascia invariata l’informazione osservativa.\end{minipage}\end{lrbox}
\MeasureBox{16-TitleBox}{\TitleBox}
\MeasureBox{16-BodyBox}{\BodyBox}
\MeasureBox{16-SideBox}{\SideBox}
\MeasureBox{16-RecallBox}{\RecallBox}
\begin{frame}[plain]\hypertarget{slide16}{}
\bookmark[level=1,page=16]{Distribuzione bootstrap della pendenza}
\begin{tikzpicture}[remember picture,overlay]\node[anchor=north east,inner sep=0] at ([xshift=19.05cm,yshift=-0.16cm]current page.north west){\color{Classical}\fontsize{8}{10}\selectfont\hyperlink{slide2}{Indice}};
\node[anchor=north west,inner sep=0] at ([xshift=0.95cm,yshift=-0.62cm]current page.north west){\usebox{\TitleBox}};
\draw[Classical,line width=1.3pt] ([xshift=.95cm,yshift=-1.68cm]current page.north west)--([xshift=2.25cm,yshift=-1.68cm]current page.north west);\node[anchor=north west,inner sep=0] at ([xshift=0.95cm,yshift=-2.12cm]current page.north west){\usebox{\BodyBox}};
\node[anchor=north west,inner sep=0] at ([xshift=11.6cm,yshift=-2.1cm]current page.north west){\color{Quantum}\bfseries\fontsize{11.5}{14}\selectfont Evidenza grafica};
\node[anchor=north west,inner sep=0] at ([xshift=11.6cm,yshift=-2.75cm]current page.north west){\usebox{\SideBox}};
\draw[Rule,line width=.5pt] ([xshift=.95cm,yshift=-10.43cm]current page.north west)--([xshift=19.05cm,yshift=-10.43cm]current page.north west);\node[anchor=north west,inner sep=0] at ([xshift=0.95cm,yshift=-10.66cm]current page.north west){\color{Classical}\bfseries\fontsize{9.4}{12}\selectfont Richiami};
\node[anchor=north west,inner sep=0] at ([xshift=2.6cm,yshift=-10.64cm]current page.north west){\usebox{\RecallBox}};
\draw[Rule,line width=.35pt] ([xshift=.95cm,yshift=-11.92cm]current page.north west)--([xshift=19.05cm,yshift=-11.92cm]current page.north west);\node[anchor=north west,inner sep=0] at ([xshift=0.95cm,yshift=-12.06cm]current page.north west){\color{Muted}\fontsize{7.7}{9.3}\selectfont P01 · Geometria dei pori e permeabilità};
\node[anchor=north east,inner sep=0] at ([xshift=17.65cm,yshift=-12.06cm]current page.north west){\color{Classical}\fontsize{7.3}{9.1}\selectfont };
\node[anchor=north east,inner sep=0] at ([xshift=19.05cm,yshift=-12.06cm]current page.north west){\color{Muted}\fontsize{7.7}{9.3}\selectfont 16/21};
\end{tikzpicture}\end{frame}
\begin{lrbox}{\TitleBox}\begin{minipage}[t]{18.2cm}\raggedright\sloppy\color{Ink}\fontsize{20.7}{24}\selectfont \bfseries\color{Navy}Confrontare intervalli costruiti diversamente\end{minipage}\end{lrbox}
\begin{lrbox}{\BodyBox}\begin{minipage}[t]{10.25cm}\raggedright\sloppy\color{Ink}\fontsize{12.4}{16.4}\selectfont L’intervallo classico comprende lo zero, mentre quello percentile resta interamente negativo. Il disaccordo rende visibile la dipendenza della conclusione da assunzioni e approssimazioni, particolarmente rilevante con dodici carote. L’analisi riporta entrambe le costruzioni e ne discute le condizioni, senza scegliere quella che produce l’esito più favorevole. Il test classico della pendenza nulla restituisce un p-value di circa 0,204, mentre il coefficiente di determinazione è circa 0,156: queste quantità descrivono aspetti differenti del modello e non eliminano il problema della limitata informazione.\end{minipage}\end{lrbox}
\begin{lrbox}{\SideBox}\begin{minipage}[t]{7.45cm}\raggedright\sloppy\color{Ink}\fontsize{12.2}{16.1}\selectfont \includegraphics[width=7.45cm]{figure/P01/F07.pdf}\par\vspace{.16cm}{\fontsize{9.5}{12.6}\selectfont Stessa stima puntuale, intervallo t e intervallo percentile. Il riferimento verticale indica la pendenza nulla.}\end{minipage}\end{lrbox}
\begin{lrbox}{\RecallBox}\begin{minipage}[t]{16.45cm}\raggedright\sloppy\color{Ink}\fontsize{10.1}{13.2}\selectfont La discordanza fra procedure orienta l’esame delle assunzioni e la richiesta di nuove evidenze.\end{minipage}\end{lrbox}
\MeasureBox{17-TitleBox}{\TitleBox}
\MeasureBox{17-BodyBox}{\BodyBox}
\MeasureBox{17-SideBox}{\SideBox}
\MeasureBox{17-RecallBox}{\RecallBox}
\begin{frame}[plain]\hypertarget{slide17}{}
\bookmark[level=1,page=17]{Confrontare intervalli costruiti diversamente}
\begin{tikzpicture}[remember picture,overlay]\node[anchor=north east,inner sep=0] at ([xshift=19.05cm,yshift=-0.16cm]current page.north west){\color{Classical}\fontsize{8}{10}\selectfont\hyperlink{slide2}{Indice}};
\node[anchor=north west,inner sep=0] at ([xshift=0.95cm,yshift=-0.62cm]current page.north west){\usebox{\TitleBox}};
\draw[Classical,line width=1.3pt] ([xshift=.95cm,yshift=-1.68cm]current page.north west)--([xshift=2.25cm,yshift=-1.68cm]current page.north west);\node[anchor=north west,inner sep=0] at ([xshift=0.95cm,yshift=-2.12cm]current page.north west){\usebox{\BodyBox}};
\node[anchor=north west,inner sep=0] at ([xshift=11.6cm,yshift=-2.1cm]current page.north west){\color{Quantum}\bfseries\fontsize{11.5}{14}\selectfont Evidenza grafica};
\node[anchor=north west,inner sep=0] at ([xshift=11.6cm,yshift=-2.75cm]current page.north west){\usebox{\SideBox}};
\draw[Rule,line width=.5pt] ([xshift=.95cm,yshift=-10.43cm]current page.north west)--([xshift=19.05cm,yshift=-10.43cm]current page.north west);\node[anchor=north west,inner sep=0] at ([xshift=0.95cm,yshift=-10.66cm]current page.north west){\color{Classical}\bfseries\fontsize{9.4}{12}\selectfont Richiami};
\node[anchor=north west,inner sep=0] at ([xshift=2.6cm,yshift=-10.64cm]current page.north west){\usebox{\RecallBox}};
\draw[Rule,line width=.35pt] ([xshift=.95cm,yshift=-11.92cm]current page.north west)--([xshift=19.05cm,yshift=-11.92cm]current page.north west);\node[anchor=north west,inner sep=0] at ([xshift=0.95cm,yshift=-12.06cm]current page.north west){\color{Muted}\fontsize{7.7}{9.3}\selectfont P01 · Geometria dei pori e permeabilità};
\node[anchor=north east,inner sep=0] at ([xshift=17.65cm,yshift=-12.06cm]current page.north west){\color{Classical}\fontsize{7.3}{9.1}\selectfont };
\node[anchor=north east,inner sep=0] at ([xshift=19.05cm,yshift=-12.06cm]current page.north west){\color{Muted}\fontsize{7.7}{9.3}\selectfont 17/21};
\end{tikzpicture}\end{frame}
\begin{lrbox}{\TitleBox}\begin{minipage}[t]{18.2cm}\raggedright\sloppy\color{Ink}\fontsize{20.7}{24}\selectfont \bfseries\color{Navy}Residui e forma della relazione\end{minipage}\end{lrbox}
\begin{lrbox}{\BodyBox}\begin{minipage}[t]{10.25cm}\raggedright\sloppy\color{Ink}\fontsize{12.4}{16.4}\selectfont I residui descrivono gli scostamenti delle risposte dalla retta adattata e vengono esaminati rispetto ai valori stimati. Curvature, variazioni di dispersione o punti isolati possono suggerire limiti della specificazione, benché dodici osservazioni offrano una capacità diagnostica ridotta. L’assenza di un andamento evidente nel grafico lascia quindi aperta l’adeguatezza delle assunzioni. I residui OLS sono ortogonali alle colonne della matrice del modello per costruzione; questa proprietà verifica il calcolo e non costituisce una prova dell’indipendenza degli errori geologici.\par\vspace{.28cm}{\centering\fontsize{13}{17}\selectfont\(\displaystyle\begin{gathered}X^\top\widehat\varepsilon=0\end{gathered}\)\par}\end{minipage}\end{lrbox}
\begin{lrbox}{\SideBox}\begin{minipage}[t]{7.45cm}\raggedright\sloppy\color{Ink}\fontsize{12.2}{16.1}\selectfont \includegraphics[width=7.45cm]{figure/P01/F05.pdf}\par\vspace{.16cm}{\fontsize{9.5}{12.6}\selectfont X: matrice con colonne 1 e percentuale areale media; $\hat\varepsilon$: vettore dei dodici residui. La linea orizzontale indica zero.}\end{minipage}\end{lrbox}
\begin{lrbox}{\RecallBox}\begin{minipage}[t]{16.45cm}\raggedright\sloppy\color{Ink}\fontsize{10.1}{13.2}\selectfont Le proprietà algebriche dell’adattamento devono essere distinte dalle ipotesi probabilistiche sugli errori.\end{minipage}\end{lrbox}
\MeasureBox{18-TitleBox}{\TitleBox}
\MeasureBox{18-BodyBox}{\BodyBox}
\MeasureBox{18-SideBox}{\SideBox}
\MeasureBox{18-RecallBox}{\RecallBox}
\begin{frame}[plain]\hypertarget{slide18}{}
\bookmark[level=1,page=18]{Residui e forma della relazione}
\begin{tikzpicture}[remember picture,overlay]\node[anchor=north east,inner sep=0] at ([xshift=19.05cm,yshift=-0.16cm]current page.north west){\color{Classical}\fontsize{8}{10}\selectfont\hyperlink{slide2}{Indice}};
\node[anchor=north west,inner sep=0] at ([xshift=0.95cm,yshift=-0.62cm]current page.north west){\usebox{\TitleBox}};
\draw[Classical,line width=1.3pt] ([xshift=.95cm,yshift=-1.68cm]current page.north west)--([xshift=2.25cm,yshift=-1.68cm]current page.north west);\node[anchor=north west,inner sep=0] at ([xshift=0.95cm,yshift=-2.12cm]current page.north west){\usebox{\BodyBox}};
\node[anchor=north west,inner sep=0] at ([xshift=11.6cm,yshift=-2.1cm]current page.north west){\color{Quantum}\bfseries\fontsize{11.5}{14}\selectfont Evidenza grafica};
\node[anchor=north west,inner sep=0] at ([xshift=11.6cm,yshift=-2.75cm]current page.north west){\usebox{\SideBox}};
\draw[Rule,line width=.5pt] ([xshift=.95cm,yshift=-10.43cm]current page.north west)--([xshift=19.05cm,yshift=-10.43cm]current page.north west);\node[anchor=north west,inner sep=0] at ([xshift=0.95cm,yshift=-10.66cm]current page.north west){\color{Classical}\bfseries\fontsize{9.4}{12}\selectfont Richiami};
\node[anchor=north west,inner sep=0] at ([xshift=2.6cm,yshift=-10.64cm]current page.north west){\usebox{\RecallBox}};
\draw[Rule,line width=.35pt] ([xshift=.95cm,yshift=-11.92cm]current page.north west)--([xshift=19.05cm,yshift=-11.92cm]current page.north west);\node[anchor=north west,inner sep=0] at ([xshift=0.95cm,yshift=-12.06cm]current page.north west){\color{Muted}\fontsize{7.7}{9.3}\selectfont P01 · Geometria dei pori e permeabilità};
\node[anchor=north east,inner sep=0] at ([xshift=17.65cm,yshift=-12.06cm]current page.north west){\color{Classical}\fontsize{7.3}{9.1}\selectfont };
\node[anchor=north east,inner sep=0] at ([xshift=19.05cm,yshift=-12.06cm]current page.north west){\color{Muted}\fontsize{7.7}{9.3}\selectfont 18/21};
\end{tikzpicture}\end{frame}
\begin{lrbox}{\TitleBox}\begin{minipage}[t]{18.2cm}\raggedright\sloppy\color{Ink}\fontsize{20.7}{24}\selectfont \bfseries\color{Navy}Escludere una carota alla volta come diagnosi\end{minipage}\end{lrbox}
\begin{lrbox}{\BodyBox}\begin{minipage}[t]{10.25cm}\raggedright\sloppy\color{Ink}\fontsize{12.4}{16.4}\selectfont La regressione viene ripetuta dodici volte, omettendo a turno una carota e mantenendo invariati modello e trasformazione. Le pendenze risultano comprese fra circa meno 0,13372 e meno 0,06707, mostrando una variazione della grandezza dell’associazione senza inversione del segno. La procedura identifica il contributo dei singoli campioni alla stima e orienta il controllo delle misure. L’esclusione definitiva richiederebbe una motivazione indipendente dalla convenienza del risultato e dovrebbe essere documentata insieme all’analisi sul campione completo.\par\vspace{.28cm}{\centering\fontsize{13}{17}\selectfont\(\displaystyle\begin{gathered}(\hat\alpha_{(-j)},\hat\beta_{(-j)})=\arg\min_{a,b}\sum_{i\ne j}(Y_i-a-b\bar A_i)^2\end{gathered}\)\par}\end{minipage}\end{lrbox}
\begin{lrbox}{\SideBox}\begin{minipage}[t]{7.45cm}\raggedright\sloppy\color{Ink}\fontsize{12.2}{16.1}\selectfont \includegraphics[width=7.45cm]{figure/P01/F06.pdf}\par\vspace{.16cm}{\fontsize{9.5}{12.6}\selectfont $\bar A_i$: percentuale areale; $Y_i=\log_{10}(K_i/1\,\mathrm{mD})$. Si omette la carota j. Viola: pendenza su tutte le dodici.}\end{minipage}\end{lrbox}
\begin{lrbox}{\RecallBox}\begin{minipage}[t]{16.45cm}\raggedright\sloppy\color{Ink}\fontsize{10.1}{13.2}\selectfont L’influenza di un’osservazione richiede esame e documentazione prima di una decisione di esclusione.\end{minipage}\end{lrbox}
\MeasureBox{19-TitleBox}{\TitleBox}
\MeasureBox{19-BodyBox}{\BodyBox}
\MeasureBox{19-SideBox}{\SideBox}
\MeasureBox{19-RecallBox}{\RecallBox}
\begin{frame}[plain]\hypertarget{slide19}{}
\bookmark[level=1,page=19]{Escludere una carota alla volta come diagnosi}
\begin{tikzpicture}[remember picture,overlay]\node[anchor=north east,inner sep=0] at ([xshift=19.05cm,yshift=-0.16cm]current page.north west){\color{Classical}\fontsize{8}{10}\selectfont\hyperlink{slide2}{Indice}};
\node[anchor=north west,inner sep=0] at ([xshift=0.95cm,yshift=-0.62cm]current page.north west){\usebox{\TitleBox}};
\draw[Classical,line width=1.3pt] ([xshift=.95cm,yshift=-1.68cm]current page.north west)--([xshift=2.25cm,yshift=-1.68cm]current page.north west);\node[anchor=north west,inner sep=0] at ([xshift=0.95cm,yshift=-2.12cm]current page.north west){\usebox{\BodyBox}};
\node[anchor=north west,inner sep=0] at ([xshift=11.6cm,yshift=-2.1cm]current page.north west){\color{Quantum}\bfseries\fontsize{11.5}{14}\selectfont Evidenza grafica};
\node[anchor=north west,inner sep=0] at ([xshift=11.6cm,yshift=-2.75cm]current page.north west){\usebox{\SideBox}};
\draw[Rule,line width=.5pt] ([xshift=.95cm,yshift=-10.43cm]current page.north west)--([xshift=19.05cm,yshift=-10.43cm]current page.north west);\node[anchor=north west,inner sep=0] at ([xshift=0.95cm,yshift=-10.66cm]current page.north west){\color{Classical}\bfseries\fontsize{9.4}{12}\selectfont Richiami};
\node[anchor=north west,inner sep=0] at ([xshift=2.6cm,yshift=-10.64cm]current page.north west){\usebox{\RecallBox}};
\draw[Rule,line width=.35pt] ([xshift=.95cm,yshift=-11.92cm]current page.north west)--([xshift=19.05cm,yshift=-11.92cm]current page.north west);\node[anchor=north west,inner sep=0] at ([xshift=0.95cm,yshift=-12.06cm]current page.north west){\color{Muted}\fontsize{7.7}{9.3}\selectfont P01 · Geometria dei pori e permeabilità};
\node[anchor=north east,inner sep=0] at ([xshift=17.65cm,yshift=-12.06cm]current page.north west){\color{Classical}\fontsize{7.3}{9.1}\selectfont };
\node[anchor=north east,inner sep=0] at ([xshift=19.05cm,yshift=-12.06cm]current page.north west){\color{Muted}\fontsize{7.7}{9.3}\selectfont 19/21};
\end{tikzpicture}\end{frame}
\begin{lrbox}{\TitleBox}\begin{minipage}[t]{18.2cm}\raggedright\sloppy\color{Ink}\fontsize{20.7}{24}\selectfont \bfseries\color{Navy}Un segno inatteso richiede una spiegazione fondata\end{minipage}\end{lrbox}
\begin{lrbox}{\BodyBox}\begin{minipage}[t]{10.25cm}\raggedright\sloppy\color{Ink}\fontsize{12.4}{16.4}\selectfont La pendenza negativa osservata non autorizza a correggere il segno delle variabili o a concludere che una maggiore porosità riduca in generale la permeabilità. La frazione areale misura una proprietà delle sezioni, mentre trasporto, connessione dei pori, tortuosità e dimensione delle gole appartengono a una struttura tridimensionale più articolata. Il campione storico e il modello lineare rappresentano ulteriori limiti. La conclusione sostiene dunque l’associazione descritta nelle carote disponibili, accompagnandola con l’incertezza e con le misure che servirebbero a esaminarne il meccanismo.\end{minipage}\end{lrbox}
\begin{lrbox}{\SideBox}\begin{minipage}[t]{7.45cm}\raggedright\sloppy\color{Ink}\fontsize{12.2}{16.1}\selectfont \input{tikz/P01/D05.tex}\end{minipage}\end{lrbox}
\begin{lrbox}{\RecallBox}\begin{minipage}[t]{16.45cm}\raggedright\sloppy\color{Ink}\fontsize{10.1}{13.2}\selectfont L’interpretazione fisica deve rispettare la proprietà effettivamente misurata e il dominio del campione.\end{minipage}\end{lrbox}
\MeasureBox{20-TitleBox}{\TitleBox}
\MeasureBox{20-BodyBox}{\BodyBox}
\MeasureBox{20-SideBox}{\SideBox}
\MeasureBox{20-RecallBox}{\RecallBox}
\begin{frame}[plain]\hypertarget{slide20}{}
\bookmark[level=0,page=20]{Argomentazione geologica e consegna}
\bookmark[level=1,page=20]{Un segno inatteso richiede una spiegazione fondata}
\begin{tikzpicture}[remember picture,overlay]\node[anchor=north east,inner sep=0] at ([xshift=19.05cm,yshift=-0.16cm]current page.north west){\color{Classical}\fontsize{8}{10}\selectfont\hyperlink{slide2}{Indice}};
\node[anchor=north west,inner sep=0] at ([xshift=0.95cm,yshift=-0.62cm]current page.north west){\usebox{\TitleBox}};
\draw[Classical,line width=1.3pt] ([xshift=.95cm,yshift=-1.68cm]current page.north west)--([xshift=2.25cm,yshift=-1.68cm]current page.north west);\node[anchor=north west,inner sep=0] at ([xshift=0.95cm,yshift=-2.12cm]current page.north west){\usebox{\BodyBox}};
\node[anchor=north west,inner sep=0] at ([xshift=11.6cm,yshift=-2.1cm]current page.north west){\color{Quantum}\bfseries\fontsize{11.5}{14}\selectfont Lettura del diagramma};
\node[anchor=north west,inner sep=0] at ([xshift=11.6cm,yshift=-2.75cm]current page.north west){\usebox{\SideBox}};
\draw[Rule,line width=.5pt] ([xshift=.95cm,yshift=-10.43cm]current page.north west)--([xshift=19.05cm,yshift=-10.43cm]current page.north west);\node[anchor=north west,inner sep=0] at ([xshift=0.95cm,yshift=-10.66cm]current page.north west){\color{Classical}\bfseries\fontsize{9.4}{12}\selectfont Richiami};
\node[anchor=north west,inner sep=0] at ([xshift=2.6cm,yshift=-10.64cm]current page.north west){\usebox{\RecallBox}};
\draw[Rule,line width=.35pt] ([xshift=.95cm,yshift=-11.92cm]current page.north west)--([xshift=19.05cm,yshift=-11.92cm]current page.north west);\node[anchor=north west,inner sep=0] at ([xshift=0.95cm,yshift=-12.06cm]current page.north west){\color{Muted}\fontsize{7.7}{9.3}\selectfont P01 · Geometria dei pori e permeabilità};
\node[anchor=north east,inner sep=0] at ([xshift=17.65cm,yshift=-12.06cm]current page.north west){\color{Classical}\fontsize{7.3}{9.1}\selectfont };
\node[anchor=north east,inner sep=0] at ([xshift=19.05cm,yshift=-12.06cm]current page.north west){\color{Muted}\fontsize{7.7}{9.3}\selectfont 20/21};
\end{tikzpicture}\end{frame}
\begin{lrbox}{\TitleBox}\begin{minipage}[t]{18.2cm}\raggedright\sloppy\color{Ink}\fontsize{20.7}{24}\selectfont \bfseries\color{Navy}Relazione, esercizio e criteri di discussione\end{minipage}\end{lrbox}
\begin{lrbox}{\BodyBox}\begin{minipage}[t]{10.25cm}\raggedright\sloppy\color{Ink}\fontsize{12.4}{16.4}\selectfont La relazione presenta provenienza, ricostruzione delle carote, trasformazione areale, modello, stima con unità, intervalli e diagnostica. L’esercizio richiede due figure motivate, il confronto fra intervallo t e bootstrap e una discussione della sensibilità alle singole carote. Una risposta adeguata collega il disaccordo inferenziale al piccolo campione e alle assunzioni, evitando di ridurlo alla scelta di una soglia di significatività. L’estensione facoltativa esamina un predittore alternativo soltanto dopo averne verificato definizione e disponibilità nei dati originali.\end{minipage}\end{lrbox}
\begin{lrbox}{\SideBox}\begin{minipage}[t]{7.45cm}\raggedright\sloppy\color{Ink}\fontsize{12.2}{16.1}\selectfont Lo script R/P01\_figure.R ricostruisce calcoli, tabelle e grafici dalla cartella dati/progetti. Le verifiche controllano gruppi, dominio, corrispondenza con rock, stima della pendenza e piano bootstrap. Il progetto raccorda misura e campionamento, inferenza e regressione, sviluppati nei moduli M01, M05 e M07.\end{minipage}\end{lrbox}
\begin{lrbox}{\RecallBox}\begin{minipage}[t]{16.45cm}\raggedright\sloppy\color{Ink}\fontsize{10.1}{13.2}\selectfont La consegna deve rendere verificabili il calcolo e il percorso che conduce dall’evidenza alla conclusione.\end{minipage}\end{lrbox}
\MeasureBox{21-TitleBox}{\TitleBox}
\MeasureBox{21-BodyBox}{\BodyBox}
\MeasureBox{21-SideBox}{\SideBox}
\MeasureBox{21-RecallBox}{\RecallBox}
\begin{frame}[plain]\hypertarget{slide21}{}
\bookmark[level=1,page=21]{Relazione, esercizio e criteri di discussione}
\begin{tikzpicture}[remember picture,overlay]\node[anchor=north east,inner sep=0] at ([xshift=19.05cm,yshift=-0.16cm]current page.north west){\color{Classical}\fontsize{8}{10}\selectfont\hyperlink{slide2}{Indice}};
\node[anchor=north west,inner sep=0] at ([xshift=0.95cm,yshift=-0.62cm]current page.north west){\usebox{\TitleBox}};
\draw[Classical,line width=1.3pt] ([xshift=.95cm,yshift=-1.68cm]current page.north west)--([xshift=2.25cm,yshift=-1.68cm]current page.north west);\node[anchor=north west,inner sep=0] at ([xshift=0.95cm,yshift=-2.12cm]current page.north west){\usebox{\BodyBox}};
\node[anchor=north west,inner sep=0] at ([xshift=11.6cm,yshift=-2.1cm]current page.north west){\color{Quantum}\bfseries\fontsize{11.5}{14}\selectfont Sviluppo e interpretazione};
\node[anchor=north west,inner sep=0] at ([xshift=11.6cm,yshift=-2.75cm]current page.north west){\usebox{\SideBox}};
\draw[Rule,line width=.5pt] ([xshift=.95cm,yshift=-10.43cm]current page.north west)--([xshift=19.05cm,yshift=-10.43cm]current page.north west);\node[anchor=north west,inner sep=0] at ([xshift=0.95cm,yshift=-10.66cm]current page.north west){\color{Classical}\bfseries\fontsize{9.4}{12}\selectfont Richiami};
\node[anchor=north west,inner sep=0] at ([xshift=2.6cm,yshift=-10.64cm]current page.north west){\usebox{\RecallBox}};
\draw[Rule,line width=.35pt] ([xshift=.95cm,yshift=-11.92cm]current page.north west)--([xshift=19.05cm,yshift=-11.92cm]current page.north west);\node[anchor=north west,inner sep=0] at ([xshift=0.95cm,yshift=-12.06cm]current page.north west){\color{Muted}\fontsize{7.7}{9.3}\selectfont P01 · Geometria dei pori e permeabilità};
\node[anchor=north east,inner sep=0] at ([xshift=17.65cm,yshift=-12.06cm]current page.north west){\color{Classical}\fontsize{7.3}{9.1}\selectfont R datasets: rock; attribuzioni e metadati in dati/progetti};
\node[anchor=north east,inner sep=0] at ([xshift=19.05cm,yshift=-12.06cm]current page.north west){\color{Muted}\fontsize{7.7}{9.3}\selectfont 21/21};
\end{tikzpicture}\end{frame}
\end{document}